\section{Justificativa}\label{sec:justificativa}
O avanço das tecnologias digitais e a popularização do ensino remoto transformaram profundamente as práticas educacionais, ampliando o acesso ao conhecimento, mas também impondo novos desafios aos professores. Entre esses desafios, destacam-se a dificuldade em manter a atenção e o engajamento dos estudantes, bem como a limitação de instrumentos que permitam observar e mensurar esses aspectos de forma sistemática.

Em ambientes virtuais de aprendizagem, a ausência do contato presencial e a variedade de estímulos externos podem favorecer a dispersão e a perda de foco, comprometendo o rendimento acadêmico e a efetividade das metodologias de ensino. Nesse contexto, torna-se importante fazer um acompanhamento ativo do discente e como ele interage nesse ambiente virtual, se faz sentido seguir com a aula expositiva nesse âmbito.

O \textit{software} ESPEON surge como uma proposta para atender a essa necessidade, monitorar o comportamento dos alunos em sala de aula virtual. A partir de métricas de tempo e estado das abas do navegador, é possível identificar padrões de engajamento e desatenção, fornecendo subsídios para aprimorar estratégias pedagógicas e compreender melhor a dinâmica da aprendizagem mediada por tecnologia.

Dessa forma, o presente estudo justifica-se pela relevância de se investigar o desempenho da metodologia expositiva no ensino remoto a partir de dados empíricos coletados em tempo real. Além de contribuir para a discussão sobre atenção e engajamento em ambientes virtuais, o trabalho propõe uma solução tecnológica de baixo custo e alta aplicabilidade, consolidando um instrumento útil para pesquisas futuras e práticas docentes.